\documentclass[11pt,a4paper]{article}

% ── Packages ──
\usepackage[utf8]{inputenc}
\usepackage[T1]{fontenc}
\usepackage{lmodern}
\usepackage[margin=2.5cm]{geometry}
\usepackage{graphicx}
\usepackage{booktabs}
\usepackage{longtable}
\usepackage{array}
\usepackage{hyperref}
\usepackage{xcolor}
\usepackage{amsmath}
\usepackage{amssymb}
\usepackage{natbib}
\usepackage{caption}
\usepackage{subcaption}
\usepackage{multirow}
\usepackage{tabularx}
\usepackage{setspace}
\usepackage{authblk}
\usepackage{lineno}

% ── npj Microgravity style approximations ──
\linenumbers
\onehalfspacing
\setlength{\parindent}{0pt}
\setlength{\parskip}{8pt}

\hypersetup{
    colorlinks=true,
    linkcolor=blue!70!black,
    citecolor=blue!70!black,
    urlcolor=blue!70!black
}

% ── Title ──
\title{Computational identification of pharmacological and nutritional countermeasures for \\
       spaceflight-induced transcriptomic changes via consensus signature reversal}

\author[1]{Richard Barker}
\author[1]{Biomni Computational Platform}
\affil[1]{Phylo, Inc.}

\date{2026-08-02}

\begin{document}
\maketitle

% ── Abstract ──
\begin{abstract}
Prolonged spaceflight induces systemic transcriptomic changes including mitochondrial dysfunction, 
cell-cycle suppression, and metabolic reprogramming. We present a reproducible computational pipeline 
(\textit{astronaut-opposite-forcing}) that constructs consensus spaceflight transcriptomic signatures 
from NASA OSDR rodent ISS RNA-seq studies via random-effects meta-analysis, screens the LINCS L1000 
perturbation compendium for compounds that reverse these signatures, annotates hits with DrugBank, 
Broad Drug Repurposing Hub, and ChEMBL, validates biological plausibility via functional enrichment, 
and translates therapeutic candidates into a vegan nutrition layer using FooDB. The liver/immune 
signature (525 genes) is dominated by downregulated E2F/G2M 
cell-cycle programs and upregulated fatty acid and xenobiotic metabolism. The brain signature 
(505 genes) shows profound downregulation of oxidative 
phosphorylation and mitochondrial complex I assembly. Top pharmacological reversers include Imatinib 
and Metformin (liver), and Vemurafenib and Formoterol (brain). Vegan nutritional countermeasures 
featuring sulforaphane, curcumin, quercetin, and EGCG provide accessible, food-based adjuncts. 
We propose a combined pharmacological and nutritional countermeasure strategy for long-duration 
spaceflight.
\end{abstract}

\noindent\textbf{Keywords:} spaceflight transcriptomics, drug repurposing, LINCS L1000, connectivity 
mapping, meta-analysis, nutritional countermeasures, vegan nutrition, mitochondrial dysfunction, 
NASA OSDR

\bigskip
\hrule
\bigskip

% ═══════════════════════════════════════════════════════
% INTRODUCTION
% ═══════════════════════════════════════════════════════
\section{Introduction}

Spaceflight imposes a unique combination of stressors on biological systems: microgravity, cosmic 
radiation, altered circadian rhythms, confinement, and modified nutrition. Multi-omics studies have 
documented widespread transcriptomic, proteomic, and metabolomic changes in astronauts and spaceflight 
model organisms \citep{garrett2019nasa, beheshti2019multi}. Key molecular hallmarks include 
mitochondrial dysfunction, oxidative stress, immune dysregulation, cell-cycle alterations, and 
metabolic reprogramming \citep{garrett2019nasa}.

The NASA Open Science Data Repository (OSDR) provides a growing collection of rodent spaceflight 
RNA-seq datasets from International Space Station (ISS) missions. While individual studies have 
identified differentially expressed genes, no consensus spaceflight transcriptomic signature has been 
built by meta-analyzing across studies. Such a consensus signature would enable systematic 
identification of countermeasures through connectivity mapping --- the principle that compounds 
producing opposite transcriptional changes may counteract the condition of interest \citep{subramanian2017next}.

The LINCS L1000 program has profiled over one million chemical and genetic perturbations, creating 
a compendium for in silico drug repurposing \citep{subramanian2017next}. The Connectivity Map (CMap) 
approach identifies compounds whose perturbation signatures reverse a disease signature --- 
``opposite forcings'' that may serve as therapeutic candidates. This approach has been successfully 
applied to numerous diseases but has not been systematically applied to spaceflight biology.

Here we present \textit{astronaut-opposite-forcing}, a reproducible pipeline that: (1) builds 
tissue-stratified consensus spaceflight signatures from rodent ISS RNA-seq studies, (2) screens LINCS 
L1000 for reversers, (3) annotates hits with multiple drug databases, (4) validates biological 
plausibility via functional enrichment, and (5) translates candidates into a vegan nutritional 
countermeasure layer. We focus on two tissue strata --- liver/immune and brain --- representing the 
metabolic and neurological vulnerabilities most relevant to long-duration spaceflight.

% ═══════════════════════════════════════════════════════
% METHODS
% ═══════════════════════════════════════════════════════
\section{Methods}

\subsection{Data acquisition}
Rodent spaceflight RNA-seq studies were identified via the NASA OSDR BDAPI 
(\url{https://osdr.nasa.gov/bio/api/}). The study catalog was filtered for \textit{Mus musculus} 
ISS missions with RNA-seq count matrices. 160 studies were screened; 31 valid STAR count matrices 
were downloaded. Human cross-check studies were cataloged for validation.

\subsection{Differential expression analysis}
Per-study differential expression was computed with DESeq2 \citep{love2014deseq2} using a 
$\sim$\texttt{condition} design with Ground Control as reference. Sample conditions were parsed from 
\texttt{sample\_metadata.json} with fallback parsing from sample names. ERCC spike-in controls were 
removed. apeglm LFC shrinkage \citep{zhu2019apeglm} was applied; studies with over-shrunk LFCs 
(median SE $< 0.01$) were excluded, retaining 15 studies with valid DE signal.

\subsection{Random-effects meta-analysis}
Gene-level random-effects meta-analysis was performed with \texttt{metafor::rma} (REML method, 
DerSimonian-Laird fallback) \citep{viechtbauer2010metafor}. Mouse-to-human ortholog mapping used 
\texttt{org.Mm.eg.db} + \texttt{homologene}. Consensus genes were defined as nominal $p < 0.05$ 
AND $\geq 70\%$ direction consistency across studies. Core genes additionally required 
$|\text{pooled log}_2\text{FC}| \geq 0.5$. Technical gene filters removed ribosomal, housekeeping, 
and translation factor genes; metallothioneins were retained.

\subsection{LINCS L1000 connectivity mapping}
Consensus signatures (human HGNC symbols) were screened against the LINCS L1000 perturbation 
compendium via the SigCom API \citep{subramanian2017next}. Gene resolution succeeded at 95--99\% 
coverage. The SigCom data API was unavailable (HTTP 503) during analysis; a local GMT fallback 
(hypergeometric overlap against \texttt{single\_drug\_perturbations-v1.0.gmt}) was used. Compounds 
were annotated with the Broad Drug Repurposing Hub (clinical phase, MoA, target, indication).

\subsection{Functional enrichment}
GSEA was performed with \texttt{clusterProfiler::GSEA} \citep{yu2012clusterprofiler} on Hallmark 
\citep{liberzon2015hallmark} and Reactome gene sets. ORA was performed with \texttt{enricher()} on 
GO:BP and Reactome. KEGG was excluded due to commercial license restrictions. Disease-gene enrichment 
used JensenLab DISEASES \citep{jensen2016diseases}. Significance thresholds: GSEA $p_{\text{adj}} < 0.25$, 
ORA $p_{\text{adj}} < 0.05$.

\subsection{Nutrient-gene mapping and vegan recipe optimization}
LINCS reverser compounds were mapped to FooDB \citep{wishart2020foodb} by compound name matching. 
Vegan food sources were classified by keyword filtering. Ten countermeasure nutrients were defined 
based on LINCS hits, FooDB coverage, and known bioactivity. Twelve curated vegan recipes were scored 
by nutrient coverage, LINCS validation, and mission-phase relevance.

% ═══════════════════════════════════════════════════════
% RESULTS
% ═══════════════════════════════════════════════════════
\section{Results}

\subsection{Consensus spaceflight signatures}
Random-effects meta-analysis of 15 rodent ISS studies yielded two tissue-stratified consensus 
signatures (Table~\ref{tab:signatures}). The liver/immune signature comprises 
525 consensus genes (283 
upregulated, 242 downregulated) with 25 
core genes. The brain signature comprises 505 consensus 
genes (208 upregulated, 297 downregulated) 
with 13 core genes. Ortholog mapping achieved 89\% coverage for 
consensus genes.

\begin{table}[ht]
\centering
\caption{Consensus spaceflight transcriptomic signatures.}
\label{tab:signatures}
\begin{tabular}{lcccc}
\toprule
\textbf{Stratum} & \textbf{Up} & \textbf{Down} & \textbf{Total} & \textbf{Core} \\
\midrule
Liver / Immune & 283 & 242 & 525 & 25 \\
Brain          & 208 & 297 & 505 & 13 \\
\bottomrule
\end{tabular}
\end{table}

\subsection{LINCS L1000 signature reversers}
Connectivity mapping identified 219 liver and 221 brain candidate reversers. Top reversers are 
shown in Table~\ref{tab:reversers}. The liver/immune signature is most strongly reversed by Imatinib 
(Bcr-Abl kinase inhibitor), Metformin (AMPK activator), and Quercetin (flavonoid antioxidant). The 
brain signature is most strongly reversed by Vemurafenib (RAF inhibitor), Formoterol ($\beta_2$-adrenergic 
agonist), and Sulforaphane (Nrf2 activator).

\begin{table}[ht]
\centering
\caption{Top LINCS L1000 signature reversers.}
\label{tab:reversers}
\small
\begin{tabular}{llll}
\toprule
\textbf{Compound} & \textbf{Z-Sum} & \textbf{Phase} & \textbf{Mechanism} \\
\midrule
\multicolumn{4}{l}{\textit{Liver / Immune}} \\
Imatinib & -3.54 & Launched & Bcr-Abl kinase inhibitor|KIT inhibi... \\
Metformin & -3.20 & Launched & insulin sensitizer \\
Glipizide & -2.87 & Launched & sulfonylurea \\
Cadmium & -2.69 & --- & --- \\
Quercetin & -2.28 & Launched & polar auxin transport inhibitor \\
\midrule
\multicolumn{4}{l}{\textit{Brain}} \\
Vemurafenib & -9.68 & Launched & RAF inhibitor \\
Chlorambucil & -7.24 & Launched & DNA inhibitor \\
Formoterol & -6.94 & Launched & adrenergic receptor agonist \\
Casiigly & -6.92 & --- & --- \\
Bicalutamide & -6.52 & Launched & androgen receptor antagonist \\
\bottomrule
\end{tabular}
\end{table}

\subsection{Functional enrichment}
GSEA Hallmark identified 32 significant liver and 23 significant brain pathways ($p_{\text{adj}} < 0.25$). 
The liver/immune signature is characterized by downregulated cell-cycle programs (E2F targets, G2M 
checkpoint, MYC targets) and upregulated metabolic pathways (fatty acid metabolism, oxidative 
phosphorylation, xenobiotic metabolism). The brain signature shows downregulated oxidative 
phosphorylation and mitochondrial complex I assembly (Table~\ref{tab:gsea}).

\begin{table}[ht]
\centering
\caption{Top GSEA Hallmark pathways.}
\label{tab:gsea}
\small
\begin{tabular}{llcc}
\toprule
\textbf{Stratum} & \textbf{Pathway} & \textbf{NES} & \textbf{$p_{\text{adj}}$} \\
\midrule
Liver & E2F TARGETS & -2.98 & 2.5e-09 \\
Liver & G2M CHECKPOINT & -2.75 & 2.5e-09 \\
Liver & FATTY ACID METABOLISM & 2.07 & 6.7e-09 \\
Liver & OXIDATIVE PHOSPHORYLATION & 2.04 & 6.7e-09 \\
Liver & XENOBIOTIC METABOLISM & 1.98 & 2.8e-08 \\
Liver & ADIPOGENESIS & 1.95 & 9.1e-08 \\
\midrule
Brain & OXIDATIVE PHOSPHORYLATION & -2.13 & 5.0e-09 \\
Brain & MYC TARGETS V1 & -1.94 & 3.6e-06 \\
Brain & MITOTIC SPINDLE & 1.88 & 3.6e-06 \\
Brain & DNA REPAIR & -1.85 & 2.7e-05 \\
Brain & MYC TARGETS V2 & -1.79 & 2.8e-03 \\
Brain & E2F TARGETS & -1.57 & 2.8e-03 \\
\bottomrule
\end{tabular}
\end{table}

\subsection{Nutrient-gene mapping and vegan countermeasures}
LINCS reverser compounds were mapped to FooDB to identify vegan food sources. Known therapeutic 
nutrients with vegan sources include Quercetin (212 vegan foods), Sulforaphane (3 foods), Ascorbic 
acid (327 foods), Luteolin (292 foods), and Curcumin (287 foods). Twelve curated vegan recipes were 
organized into pre-flight, in-flight, and post-flight meal plans. The top-scoring recipe (Broccoli 
Sulforaphane Power Bowl, score=143) targets sulforaphane, curcumin, ascorbic acid, quercetin, and 
lutein.

% ═══════════════════════════════════════════════════════
% DISCUSSION
% ═══════════════════════════════════════════════════════
\section{Discussion}

The consensus signatures recapitulate known spaceflight biology with tissue specificity. The 
liver/immune signature's downregulation of E2F targets and G2M checkpoint genes reflects cell-cycle 
suppression observed across multiple spaceflight studies. The concurrent upregulation of fatty acid 
metabolism and oxidative phosphorylation suggests a hepatic metabolic shift toward energy storage 
and detoxification.

The brain signature's downregulation of oxidative phosphorylation and mitochondrial complex I assembly 
is particularly significant. Mitochondrial dysfunction is one of the most consistently reported 
molecular effects of spaceflight, observed in rodent studies, the NASA Twin Study, and in vitro 
experiments \citep{garrett2019nasa}. The downregulation of complex I specifically implicates impaired 
electron transport chain function.

The LINCS reversers span diverse mechanisms. Imatinib's reversal of the liver signature may relate 
to PDGFR signaling in hepatic stellate cells. Metformin activates AMPK, counteracting metabolic 
reprogramming and providing mitochondrial protection. Formoterol, a $\beta_2$-adrenergic agonist, 
may counteract muscle atrophy signaling and has neuroprotective properties. Sulforaphane activates 
Nrf2 antioxidant response, directly addressing oxidative stress.

The vegan nutritional countermeasures provide an accessible, low-risk adjunct to pharmacological 
intervention. Sulforaphane, curcumin, quercetin, and EGCG all have documented anti-inflammatory, 
antioxidant, and mitochondrial protective effects. The phase-based recipe design aligns nutritional 
intervention with the temporal dynamics of spaceflight adaptation.

\subsection{Limitations}
The LINCS screening used a local GMT fallback rather than the full SigCom API, which was unavailable. 
Meta-analysis heterogeneity ($I^2 \approx 99\%$) is high, reflecting pooling across different tissues, 
missions, and durations. Ortholog mapping achieved 89\% coverage for consensus genes but only 39.1\% 
genome-wide. JensenLab DISEASES enrichment yielded limited results because spaceflight is not a 
disease ontology term. RecipeNLG was not programmatically accessible; curated recipes were used instead.

% ═══════════════════════════════════════════════════════
% CONCLUSIONS
% ═══════════════════════════════════════════════════════
\section{Conclusions}

This pipeline demonstrates a reproducible computational approach to identifying countermeasures for 
spaceflight-induced biological changes. The consensus signatures capture tissue-specific 
vulnerabilities --- hepatic metabolic reprogramming and cerebral mitochondrial dysfunction --- 
consistent with the broader spaceflight biology literature. The combination of pharmacological 
reversers and vegan nutritional countermeasures suggests a multi-modal countermeasure strategy for 
long-duration spaceflight. Future work should validate top reversers in spaceflight-analog models, 
expand the human cross-check, and refine recipe optimization with bioavailability and shelf-stability 
constraints.

% ═══════════════════════════════════════════════════════
% REFERENCES
% ═══════════════════════════════════════════════════════
\begin{thebibliography}{99}

\bibitem{garrett2019nasa}
Garrett-Bakelman, F.E.\ et al. The NASA Twins Study: A multidimensional analysis of a year-long 
human spaceflight. \textit{Science} \textbf{364}, 6436 (2019).

\bibitem{beheshti2019multi}
Beheshti, A.\ et al. Multi-omics and systems biology approaches for understanding spaceflight 
effects. \textit{Cell} (2019).

\bibitem{subramanian2017next}
Subramanian, A.\ et al. A next generation connectivity map: L1000 platform and the first 
1,000,000 profiles. \textit{Cell} \textbf{171}, 1437--1452 (2017).

\bibitem{love2014deseq2}
Love, M.I., Huber, W.\ \& Anders, S. Moderated estimation of fold change and dispersion for 
RNA-seq data with DESeq2. \textit{Genome Biology} \textbf{15}, 550 (2014).

\bibitem{zhu2019apeglm}
Zhu, Q.\ et al. apeglm: approximate posterior estimation for RNA-seq log fold changes. 
\textit{Nature Methods} (2019).

\bibitem{viechtbauer2010metafor}
Viechtbauer, W. Conducting meta-analyses in R with the metafor package. \textit{Journal of 
Statistical Software} \textbf{36}, 1--48 (2010).

\bibitem{liberzon2015hallmark}
Liberzon, A.\ et al. The Molecular Signatures Database (MSigDB) Hallmark Gene Set Collection. 
\textit{Cell Systems} \textbf{1}, 417--425 (2015).

\bibitem{yu2012clusterprofiler}
Yu, G.\ et al. clusterProfiler: an R package for comparing biological themes among gene clusters. 
\textit{OMICS} \textbf{16}, 284--287 (2012).

\bibitem{jensen2016diseases}
Jensen, L.J.\ et al. JensenLab DISEASES: text mining and data integration for disease-gene 
associations. \textit{Database} (2016).

\bibitem{wishart2020foodb}
Wishart, D.S.\ et al. FooDB: a food compound database. \url{https://foodb.ca} (2020).

\end{thebibliography}

\end{document}
